\documentclass[10pt,a4paper]{amsart}
\usepackage[margin=19mm]{geometry}
\usepackage{amsmath,amssymb,mathtools}
\usepackage[scr=rsfs,cal=euler]{mathalfa}
\usepackage{xfrac}
\usepackage[unicode,hidelinks,bookmarks=false]{hyperref}
\newcommand{\ie}{i.e.\ }
\newcommand{\cf}{cf.\ }
\newcommand{\resp}{respectively\ }
\newcommand{\Subsection}{\subsection}
\providecommand{\nobreakdash}{\nobreak\hbox{-}}
\setlength{\emergencystretch}{2em}
\multlinegap0pt
\usepackage{bm}
\usepackage{bbm}
\usepackage{dsfont}
\usepackage{xspace}
\usepackage{cancel}
\usepackage{mathtools}
\usepackage{enumitem}
\usepackage{amsthm}
\makeatletter
\@ifundefined{theorem}{\newtheorem{theorem}{Theorem}}{}
\@ifundefined{prop}{\newtheorem{prop}[theorem]{Proposition}}{}
\makeatother
\newcommand{\X}{\times}
\newcommand{\Zd}{\mathds{Z}^4_\uparrow}
\newcommand{\Z}{\mathds{Z}}
\newcommand{\al}{\alpha}
\newcommand{\be}{\beta}
\newcommand{\bx}{\mathbf{x}}
\newcommand{\by}{\mathbf{y}}
\newcommand{\id}{\text{id}}
\newcommand{\lf}{\left}
\newcommand{\mathand}{\;\text{and}\;}
\newcommand{\mathfor}{\;\text{for}\;}
\newcommand{\rg}{\right}
\newcommand{\scrF}{\mathcal{F}}
\newcommand{\scrH}{\mathcal{H}}
\newcommand{\scrN}{\mathcal{N}}
\newcommand{\sig}{\sigma}
\newcommand{\sset}{\subset}
\title[Hidden invariance of last passage percolation and directed p]{Hidden invariance of last passage percolation and directed polymers}
\author{Duncan Dauvergne}
\date{Source: arXiv:2002.09459v5 (CC BY 4.0)}
\begin{document}
\maketitle

\noindent\emph{Source and licence.} Hidden invariance of last passage percolation and directed polymers. Version arXiv:2002.09459v5 (CC BY 4.0), distributed under the Creative Commons Attribution 4.0 International licence, as recorded in the arXiv licence metadata for this exact version and shown on the abstract page https://arxiv.org/abs/2002.09459v5. Full rights notice: licenses/arxiv-2002.09459-CC-BY-4.0.txt.\par\smallskip

\noindent\textit{Editorial scope.} Counted unit: the Prop whose source label is \texttt{P:lie-in-F} in the article (source0.tex lines 1515--1521, source label \texttt{P:lie-in-F}), delivered together with the article's own complete original proof of it (source0.tex lines 1523--1566, one \texttt{proof} environment of 44 lines). No mathematics is written, repaired or re-derived: statement and proof are the article's own text line for line. Required context retained uncounted: P:basic-perm (lines 1440--1446); P:slides (lines 1470--1473). Every definition, display and label of those lines is retained unchanged. Omitted from this package and not redistributed: 1--1439, 1447--1469, 1474--1514, 1522--1522, 1567--2070. Nothing inside a retained excerpt is omitted or altered: every retained body is delivered exactly as the article's lines are written, so no omission receipt of any kind accompanies this package. Citations: the retained excerpts carry no citation command at all, so no bibliography entry of the article is transcribed and no reference is redistributed. Reference adaptation: none was needed and none was made; every reference inside the delivered unit resolves inside it, so no unresolved reference or citation is produced. Printed slips kept literal: no printed word, symbol, hypothesis or number of the counted statement or of its proof was altered. Layout and class: the article's own class is \texttt{imsart} with packages including amsmath, amsthm, amssymb, enumitem, placeins, graphicx, float, caption, subcaption, epstopdf, inputenc, babel, dsfont, natbib; the wrapper is the lane's pinned \texttt{amsart} editorial wrapper at 10pt on A4, which restates the theorem-like environments the retained text needs and adds the article's own macro definitions that the retained text needs (\texttt{\textbackslash X}, \texttt{\textbackslash Z}, \texttt{\textbackslash Zd}, \texttt{\textbackslash al}, \texttt{\textbackslash be}, \texttt{\textbackslash bx}, \texttt{\textbackslash by}, \texttt{\textbackslash id}, \texttt{\textbackslash lf}, \texttt{\textbackslash mathand}, \texttt{\textbackslash mathfor}, \texttt{\textbackslash rg}, \texttt{\textbackslash scrF}, \texttt{\textbackslash scrH}, \texttt{\textbackslash scrN}, \texttt{\textbackslash sig}, \texttt{\textbackslash sset}), with extra wrapper packages (bm, bbm, dsfont, xspace, cancel, mathtools, enumitem). The delivered numbering is the unit's own. Assets: no retained span contains a figure environment, a graphics-inclusion command, a PDF-page inclusion, a table or a TikZ picture; no image file is copied into this package, the package declares no asset dependency and no image file is redistributed with it. Licence: version arXiv:2002.09459v5 of this article is distributed under the Creative Commons Attribution 4.0 International licence, as recorded in the arXiv licence metadata for this exact version and shown on the abstract page https://arxiv.org/abs/2002.09459v5. A full rights notice is delivered at licenses/arxiv-2002.09459-CC-BY-4.0.txt. Characters: every retained excerpt is pure ASCII, so the delivered engine, XeLaTeX with its default Latin Modern text font, reports no missing character.
\begin{prop}
	\label{P:basic-perm} (Column Transposition)
	Let
	$U = U_h \cup U_n, W = W_a \cup W_b$ be subsets of $\Zd$. Suppose that $W_a \X W_b, U_n \X W \sset \scrN,$ and $U_h \X W \sset \scrH$. Suppose also that the sets $W_a$ and $W_b$ are connected, and that there is some box $b \in \Zd$ such that $\sqcup W \sset b$ and $\sqcup U_n \sset b^c$.
	
	Let $f$ be a map with domain $U \cup W$ such that $f|_U = \id, f|_{W_a} = T_{(k, 0)},$ and $f|_{W_b} = T_{(\ell, 0)}$ for some $k, \ell \in \Z$. Suppose that $f$ preserves $\scrH$ and $\scrN$, and that $\sqcup f(W) \sset b$. Then $f \in \scrF$.
\end{prop}

\begin{prop}
	\label{P:slides} (Slides)
	Let $U, W \sset \Zd$. Suppose that we can partition $U = U_1 \cup U_2$ and $W = W_1 \cup W_2$ such that $U_1 \X W_1, W_2 \X U_2 \sset \scrH \cup \scrN$ and $U_1 \X W_2, U_2 \X W_1 \sset \scrN$. Now let $c \in \{(0, \pm 1), (\pm 1, 0\}$, and define a map $\sig$ with domain $U \cup W$ by  $\sig|_U = T_c, \sig|_W = \id$. Suppose that $\sig$ preserves $\scrH$ and $\scrN$. Then $\sig \in \scrF$.
\end{prop}

\begin{prop}
	\label{P:lie-in-F}
	\begin{enumerate}
		\item (Towers) Let $U_1, \dots, U_k, U_1', \dots, U_k' \sset \Zd$ be such that $U_i \X U_{i+1}, U'_i \X U'_{i+1} \sset \scrH$ for all $i$, and such that for all $i$, $U_i' = T_{c_i} U_i$ for some $c_i \in \Z^2$. Let $f:\bigcup U_i \to \bigcup U_i'$ be the function whose restriction to $U_i$ is $T_{c_i}$. Then $f \in \scrF$. 
		\item (Box permutations) Let $U_1, \dots, U_k, V_1, \dots, V_m, U_1', \dots, U_k', V_1', \dots, V_m' \sset \Zd$ be such that $U_i \X V_j, U_i' \X V_j' \sset \scrH$ for any $i, j$, and $U_i \X U_j, V_i \X V_j, U_i' \X U_j', V_i' \X V_j' \sset \scrN$ for any $i \ne j$. Suppose also that for all $i, j$, there exists $c_i, d_j \in \Z^2$ such that $U_i' = T_{c_i} U_i$ and $V_j' = T_{d_j} V_j$. Let $f:\bigcup U_i \cup \bigcup V_i \to \bigcup U_i' \cup \bigcup V_i'$ be the function with $f|_{U_i} = T_{c_i}$ and $f|_{V_j} = T_{d_j}$. Then $f \in \scrF$.
	\end{enumerate}
\end{prop}

\begin{proof}
	\textbf{Proof of 1.} (Towers):  \qquad We proceed inductively starting with $k = 2$ (the $k=1$ case is trivial). First, by translating $U_1', U_2'$ by a common amount, we may assume that $c_2 = 0$. Second, when $c_1 \in \{(\pm 1, 0), (0, \pm 1)\}$, then the fact that $f \in \scrF$ follows from Proposition \ref{P:slides}. For the case of general $c_1$, observe that since $U_1 \X U_2, T_{c_1} U_1 \X U_2 \sset \scrH$, that $T_x U_1 \X U_2 \sset \scrH$ for any $x$ in the box with two diagonally opposite corners given by $c_1$ and $(0, 0)$. Therefore we can compose $f$ of slide maps $f_i:T_{x_i} U_1 \X U_2 \to T_{x_{i+1}} U_1 \X U_2$ where $x_0, \dots, x_k$ is a path from $(0, 0)$ to $c_1$ in that box.
	
	For the inductive step when $k \ge 3$, we can create the map $f:\bigcup U_i \to \bigcup U_i'$ by composing two maps of the form in (1), where the first has $c_1 = c_2$, and the second has $c_2 = c_3 = \dots = c_k$. Each of these maps lies in $\scrF$ by the inductive hypothesis.
	
	\textbf{Proof of 2.} (Box permutations): \qquad  We may assume that each of the sets $U_i, V_i$ is connected. If not, we can simply further break them up into their connected components and prove that the corresponding larger set of functions lies in $\scrF$. 
	
	Fix a box $v = (v^-, v^+) \in \bigcup V_j$. For every $x \in [v^-_2, v^+_2]$, there is at most one value of $i$ for which $[v^-_1, v^+_1] \X \{x\}$ intersects $\sqcup U_i$. Moreover, for each $i$, the connectivity of each of the sets $U_i$ and the fact that $U_i \X \{v\} \sset \scrH$ for all $i$ implies that the set of such $x$ for which  $[v^-_1, v^+_1] \X \{x\}$ intersects $\sqcup U_i$ is an interval $I_i = [I_i^-, I_i^+]$ given by the projection of $\sqcup U_i$ onto the second coordinate. In particular, $I_i$ is independent of the choice of $v$.
	
	We similarly define intervals $J_i = [J_i^-, J_i^+]$ given by projecting $\sqcup V_i$ onto the first coordinate, and intervals $I_i'$ and $J_i'$ corresponding to $U_i'$ and $V_i'$. By possibly relabelling, we may assume that $I_i^+ < I^-_{i+1}$ and $J_i^+ < J^-_{i+1}$ for all $i$.
	
	\textbf{A special case:} $I_i = I_i', J_i = J_j'$ for all $i, j$. \qquad Let $T_i$ be the largest interval such that every $u \in U_i$ crosses $T_i \X I_i$ horizontally and let $S_i$ be the largest interval such that every $v \in V_i$ crosses $J_i \X S_i$ vertically. Similarly define $T_i', S_i'$. Observe that the intervals $T_i, I_i$ determine the location of $U_i$ among the set of all possible translations. Note that $T_i' = \al_i + T_i, S_i' = \be_i + S_i$ for integers $\al_i, \be_i$, and since $I_i = I_i', J_i = J_i'$, we have
	$$
	U_i' = T_{(\al_i, 0)} U_i \qquad \mathand \qquad  V_i' = T_{(0, \be_i)} V_i.
	$$
	Now let $S_1 = \prod_{i=1}^k [0, \al_i], S_2 =  \prod_{i=1}^m [0, \be_i],$ where the order of the endpoints are switched whenever $\be_i, \al_i < 0$.
	For any $(\bx, \by) \in S_1 \X S_2$, observe that $T_{(x_i, 0)} U_i \X T_{(0, y_j)} V_j \sset \scrH$ for any $i, j$, and that 
	$$
	T_{(x_i, 0)} U_i \X T_{(x_j, 0)} U_j, \quad T_{(0, y_i)} V_i \X T_{(0, y_j)} V_j \sset \scrN \quad \mathfor i \ne j.
	$$ 
	Therefore 
	for any $(\bx, \by), (\bx', \by') \in S_1 \X S_2$ with $|(\bx, \by) - (\bx', \by')| = 1$, the slide map
	$$
	g:\bigcup T_{(x_i, 0)} U_i \cup \bigcup T_{(0, y_i)} V_i \to \bigcup T_{(x_i', 0)} U_i \cup \bigcup T_{(0, y_i')} V_i
	$$
	which is the identity everywhere except for at the single coordinate where $(\bx, \by)_i \ne (\bx', \by')_i$, lies in $\scrF$ by Proposition \ref{P:slides}. We can compose the map $f:\bigcup U_i \cup \bigcup V_i \to \bigcup U_i' \cup \bigcup V_i'$ from such maps.
	
	\textbf{The general case:} \qquad By the first case, we can first apply a transformation to get that $T_i = [J_1^-, t_i]$ and $S_j = [I_1^-, s_i]$ for all $i, j$. By translating, we may also assume that 
	$$
	I_1^- = \min \lf\{x \in \Z: x \in \bigcup I_i' \rg\}, \quad \mathand \quad J_1^- = \min \lf\{x \in \Z: x \in \bigcup J_i' \rg\}.
	$$
	By again applying the first case we may then assume that $T'_i = T_i, S_i = S_i'$.
	
	Now, given $S_i, T_i$ as above, letting $\hat s = \min_i s_i$ and $\hat t = \min_i t_i$, the only constraints on the intervals $I_i'$ are that they are disjoint and contained in $[I_1^-, \hat s]$ and the only constraints on the $J_i'$ are that they are disjoint and contained in $[J_1^-, \hat t]$. We can generate all such collections of intervals from $I_i, J_i$ if we can change the order of any two adjacent intervals and translate any interval without changing the order of the intervals. 
	
	We check that such moves are in $\scrF$. By symmetry, it suffices to check that we can move the intervals $J_i$. Fix $i \in \{1, \dots m-1\}$ and let $J \sset [J_1^-, \hat t]$ be the largest interval containing $J_i \cup J_{i+1}$ that is disjoint from all the other $J_j$.
	Then any map $f$ with domain $\bigcup U_i \cup \bigcup V_i$ such that:
	\begin{itemize}[nosep]
		\item $f|_{V_i} = T_{(\ell, 0)}, f|_{V_{i+1}} = T_{(k, 0)}$, and $f$ is the identity everywhere else,
		\item $\sqcup (f(V_1) \cup f(V_2)) \sset J \X \Z$,
		\item $f$ preserves $\scrH$ and $\scrN$,
	\end{itemize}
	is in $\scrF$ by Proposition \ref{P:basic-perm}. In the application of the proposition, we take $\bigcup U_i = U_h, \bigcup_{j \ne i, i+1} V_j = U_n, V_i = W_a,$ and $V_{i+1} = W_b$. The box $b$ in the proposition can be any large enough box of the form $J \X [-r, r]$. Such maps $f$ allow us to change the order of any two intervals and translate any interval without changing the interval order, as desired.
\end{proof}

\end{document}
